\documentclass[12pt, aspectratio=169]{beamer}
\usetheme{CambridgeUS}
\usecolortheme{beaver}
\setbeamertemplate{navigation symbols}{}
\usepackage[utf8]{inputenc}
\usepackage{graphicx}
\usepackage{booktabs}

\title[SLR on Cyber Deception]{A Systematic Literature Review on Cyber Deception}
\subtitle{From Game Theory to Generative AI}
\author{Alexandru Nicolescu}
\institute[Unibo]{University of Bologna, MSc in Computer Science}
\date{February 2026}

\begin{document}
\maketitle

\begin{frame}{Agenda}
    \tableofcontents
\end{frame}

\begin{frame}{Research Methodology}
    \begin{columns}
        \column{0.5\textwidth}
        \textbf{Search Strategy}
        \begin{itemize}
            \item \textbf{Sources:} Google Scholar, IEEE Xplore, ACM DL, USENIX.
            \item \textbf{Keywords:} "Cyber Deception", "Honeypots", "MTD", combined with "LLM", "Game Theory", "RL".
            \item \textbf{Time Window:} 2002 (Foundational) -- 2025 (Emerging).
        \end{itemize}

        \column{0.5\textwidth}
        \textbf{Quality Assessment (Tiered Approach)}
        To balance authority and novelty:
        \begin{itemize}
            \item \textbf{Tier 1 (Gold Standard):} Top-tier venues (CORE A*/A) like USENIX Security, IEEE TIFS, ACM CSUR.
            \item \textbf{Tier 2 (Emerging):} Pre-prints (arXiv) for cutting-edge LLM/AI research (2023-2025).
            \item \textbf{Tier 3:} Contextual contributions.
        \end{itemize}
    \end{columns}
\end{frame}

\section{Introduction}
\begin{frame}{Context: The Need for a Third Line of Defense}
    \begin{itemize}
        \item \textbf{Current Landscape:} A "State of War". Traditional defenses (Access Control, IDS) are failing against modern threats.
        \item \textbf{The Gap:} IDS are reactive and signature-based. They fail against zero-days.
        \item \textbf{Cyber Deception:} A proactive "Active Defense" strategy.
              \begin{block}{Definition}
                  Deliberate manipulation of the attacker's perception, beliefs, or actions to cause failure or resource waste.
              \end{block}
        \item \textbf{Evolution:} From simple, static honeypots to complex, AI-driven adaptive ecosystems.
    \end{itemize}
\end{frame}

\section{Taxonomy of Deception Tactics}
\begin{frame}{Taxonomy: Conceptual Frameworks}
    Deception can be categorized by \textit{mechanism} (How it works):

    \vspace{0.5cm}
    \begin{columns}
        \column{0.5\textwidth}
        \textbf{1. Crypsis (Hiding the Real)}
        \begin{itemize}
            \item \textbf{Perturbation:} Adding noise to data.
            \item \textbf{MTD:} Shifting attack surface.
            \item \textbf{Obfuscation:} Hiding routing/assets.
        \end{itemize}

        \column{0.5\textwidth}
        \textbf{2. Mimesis (Showing the False)}
        \begin{itemize}
            \item \textbf{Static (Honey-X):} Decoys, Honey-tokens.
            \item \textbf{Dynamic (Engagement):} Interactive story-telling to guide attackers.
        \end{itemize}
    \end{columns}
\end{frame}

\begin{frame}{Taxonomy: Deployment Layers}
    Deception operates across three distinct layers:

    \begin{enumerate}
        \item \textbf{Network \& Protocol Level:}
              \begin{itemize}
                  \item Topology Manipulation (MTD), Protocol Emulation (Honeyd).
                  \item fake services via ARP spoofing.
              \end{itemize}

        \item \textbf{Host \& Application Level:}
              \begin{itemize}
                  \item Server-side Honeypots, Honey-tokens (fake AWS keys).
                  \item Shadow Honeypots (redirecting traffic).
              \end{itemize}

        \item \textbf{Cognitive \& Semantic Level:}
              \begin{itemize}
                  \item Targeting the \textit{human} mind.
                  \item Fake errors, linguistic deception to induce hesitation.
              \end{itemize}
    \end{enumerate}
\end{frame}

\section{Reactive vs. Proactive Strategies}
\begin{frame}{Strategic Evolution: Reactive vs. Proactive}
    \begin{columns}
        \column{0.5\textwidth}
        \begin{alertblock}{Reactive (Traditional)}
            \begin{itemize}
                \item \textbf{Model:} Static, passive traps.
                \item \textbf{Goal:} Forensics \& Detection.
                \item \textbf{Limit:} Cedes initiative to the attacker. High fingerprinting risk.
            \end{itemize}
        \end{alertblock}

        \column{0.5\textwidth}
        \begin{exampleblock}{Proactive (Modern)}
            \begin{itemize}
                \item \textbf{Model:} Anticipates \& Manipulates.
                \item \textbf{Goal:} Deterrence \& Misdirection.
                \item \textbf{Enablers:} Game Theory (Signaling), MTD, Cognitive biases.
            \end{itemize}
        \end{exampleblock}
    \end{columns}
    \vspace{0.5cm}
    \centering
    \textit{The shift involves moving from "detecting the breach" to "managing the adversary".}
\end{frame}

\section{Deception Placement Timing Models}
\begin{frame}{Timing Models: When to Deploy?}
    The placement of deceptive artifacts has evolved from design-time to run-time.

    \begin{itemize}
        \item \textbf{Static (Design-Time):}
              \begin{itemize}
                  \item Manual configuration based on assumptions.
                  \item Vulnerable to reconnaissance and "set it and forget it" failure.
              \end{itemize}

        \item \textbf{Dynamic (Run-Time):}
              \begin{itemize}
                  \item \textbf{On-Demand Instantiation:} Virtual honeypots created instantly upon scan (e.g., Honeyd).
                  \item \textbf{Game-Theoretic Optimization:} Calculating optimal placement based on live attacker moves.
                  \item \textbf{Adaptive Deployment (RL):} Autonomous agents learning optimal defense policies in real-time.
              \end{itemize}
    \end{itemize}
\end{frame}

\section{Attacker Profiling \& Cognitive Modeling}
\begin{frame}{Targeting the Adversary}
    Modern systems move beyond "one-size-fits-all" traps to model the attacker.

    \vspace{0.5cm}
    \textbf{1. Explicit Cognitive Modeling:}
    \begin{itemize}
        \item Disrupting the \textbf{OODA Loop} (Observe-Orient-Decide-Act).
        \item Exploiting biases: \textit{Confirmation Bias} (planting data that confirms false hypotheses) and \textit{Decoy Effect}.
    \end{itemize}

    \textbf{2. Implicit Modeling (Game Theory \& AI):}
    \begin{itemize}
        \item \textbf{Signaling Games:} Manipulating attacker's belief distribution via "Cheap Talk" (costless deceptive signals).
        \item \textbf{RL Agents:} Implicitly learning attacker profiles by observing state transitions and reward maximization.
    \end{itemize}
\end{frame}

\begin{frame}{Hacking the Human: Cognitive Deception}
    Beyond algorithms, deception targets the human operator's \textbf{OODA Loop} (Observe-Orient-Decide-Act).

    \begin{block}{Exploiting Cognitive Biases}
        \begin{itemize}
            \item \textbf{Confirmation Bias:} Planting false data that confirms the attacker's pre-existing hypothesis, reinforcing their wrong path.
            \item \textbf{Decoy Effect:} Presenting a clearly inferior option to manipulate the choice between other targets.
        \end{itemize}
    \end{block}

    \textbf{Goal-Oriented Deception:}
    Understanding the attacker's specific objective (e.g., "steal database") and constructing a high-effort, fabricated path towards it.
\end{frame}

\begin{frame}{The Mechanics of Deception: Signaling Games}
    Deception is modeled as a \textbf{Signaling Game} with incomplete information.

    \vspace{0.5cm}
    \begin{itemize}
        \item \textbf{The Players:}
              \begin{itemize}
                  \item \textbf{Sender (Defender):} Possesses private information (True Network State).
                  \item \textbf{Receiver (Attacker):} Updates beliefs based on observed signals.
              \end{itemize}

        \item \textbf{The Strategy ("Cheap Talk"):}
              The defender transmits costless, misleading signals (e.g., fake traffic patterns) to alter the attacker's belief distribution.

        \item \textbf{Goal:} Maximize the probability that the attacker chooses a suboptimal action based on falsified intelligence.
    \end{itemize}
\end{frame}

\section{Role of AI, LLMs, and RL}
\begin{frame}{AI for Defense: Adaptive \& Generative}
    \textbf{Deep Reinforcement Learning (DRL):}
    \begin{itemize}
        \item Frameworks use DQN/PPO to autonomously orchestrate defense.
        \item Cycle: Data Ingestion $\rightarrow$ DRL Decision $\rightarrow$ Response Execution.
    \end{itemize}

    \textbf{Generative Honeypots (LLMs):}
    \begin{itemize}
        \item \textit{HoneyGPT, shelLM}: Use LLMs to generate realistic terminal responses dynamically.
        \item Overcomes the "static signature" problem of traditional honeypots.
        \item \textbf{Limitations:} Context window "amnesia" (forgetting file creation) and hallucinations (fake syntax errors).
    \end{itemize}
\end{frame}

\begin{frame}{AI as the Threat: The Offensive Arms Race}
    Defenders must adapt to AI-driven adversaries:

    \begin{itemize}
        \item \textbf{Automated Pentesting:} Tools like \textit{PentestGPT} maintain context and automate complex attack chains (228\% improvement over GPT-4).
        \item \textbf{Autonomous Agents:} LLM agents utilizing ReAct frameworks to exploit 1-day vulnerabilities autonomously.
    \end{itemize}

    \vspace{0.5cm}
    \begin{alertblock}{Implication}
        Future deception must be designed not just to fool humans, but to poison the reasoning logic of AI agents.
    \end{alertblock}
\end{frame}

\section{Future Challenges}
\begin{frame}{Open Challenges \& Research Directions}
    \textbf{1. Predictive \& Intent-Aware Deception:}
    \begin{itemize}
        \item Moving from reactive RL to truly predictive models that anticipate intent \textit{before} reconnaissance.
    \end{itemize}

    \textbf{2. Domain-Specific Gaps (OT/ICS):}
    \begin{itemize}
        \item Lack of high-fidelity deception for Critical Infrastructure.
        \item \textit{Fingerprinting:} Generic honeypots leak IT-stack artifacts (e.g., Window Scaling) inconsistent with real OT hardware.
    \end{itemize}

    \textbf{3. Metrics \& Evaluation:}
    \begin{itemize}
        \item Absence of standardized metrics (TTPs, Time-to-Compromise) makes comparing strategies impossible.
    \end{itemize}
\end{frame}

\begin{frame}{The Evaluation Crisis: Lack of Standard Metrics}
    A major barrier to adoption is the inability to compare deception strategies.

    \begin{itemize}
        \item \textbf{Current State:} Studies use heterogeneous, ad-hoc criteria (e.g., "number of packets captured", "malware samples").
        \item \textbf{The Missing Standard:} We lack unified metrics for:
              \begin{itemize}
                  \item \textbf{Attacker Cost:} Time/Resources wasted.
                  \item \textbf{Believability:} Probability of deception detection.
                  \item \textbf{Enticement:} Ratio of attackers lured vs. total threats.
              \end{itemize}
        \item \textbf{Consequence:} Impossible to answer basic questions like "Is strategy A more cost-effective than strategy B?".
    \end{itemize}
\end{frame}

\section{Conclusion}
\begin{frame}{Conclusion}
    \begin{enumerate}
        \item \textbf{Evolution:} Cyber Deception has matured from static traps to dynamic, intelligent ecosystems.
        \item \textbf{Paradigm Shift:} The focus has moved from \textit{detection} to \textit{active manipulation} of the adversary's reality.
        \item \textbf{The Future:} The battleground is shifting to \textbf{AI vs. AI}. Defensive systems must become autonomous, predictive, and high-fidelity to counter machine-speed attacks.
    \end{enumerate}

\end{frame}

\begin{frame}[plain]
    \centering
    \Huge \textbf{Thank you for your attention}

    \vspace{4cm}
    \normalsize
    Alexandru Nicolescu \\
    \textit{alexandru.nicolescu@studio.unibo.it}
\end{frame}
\end{document}